% SGA 6，第 I 讲，PDF 122—149 页，第 3—5 节。
\SourcePage{122}
\SGAsection{3}{与经典凝聚性概念的联系}
\textbf{3.0。}假设与记号同（2.0），但不再假定 $\mathcal C_0$ \Term{local-scheme}{局部}地对\Term{epimorphism}{满态射}的核\Term{quasi-stable}{拟稳定}。

\begin{SGAresult}{定义}{3.1}
设 $F\in\operatorname{ob}\mathcal C_S$。若对 $S$ 上方的每个 $U$ 以及每个\Term{morphism}{态射} $u:E\to F|U$，其中 $E\in\operatorname{ob}\mathcal C_{0U}$，$\operatorname{Ker}(u)$ 都是 $\mathcal C_0$-\Term{finite-type}{有限型}的（1.2），则称 $F$ 为\Term{precoherent}{预凝聚}的（相对于 $\mathcal C_0$）。\footnote{这里引入此概念仅用作技术上的中间步骤；真正需要保留的概念只有\Term{coherent}{凝聚性}。}若 $F$ 是\Term{precoherent}{预凝聚}且\Term{finite-type}{有限型}的，则称 $F$ 为\Term{coherent}{凝聚}的（相对于 $\mathcal C_0$）。
\end{SGAresult}

\textbf{注记 3.2。}设 $F\in\operatorname{ob}\mathcal C_S$。由定义显然可知，若 $F$ 为\Term{precoherent}{预凝聚}的（相应地，\Term{finite-type}{有限型}的），则 $F$ 的每个\Term{subobject}{子对象}（相应地，\Term{quotient-object}{商对象}）也为\Term{precoherent}{预凝聚}的（相应地，\Term{finite-type}{有限型}的）。

\begin{SGAresult}{\Term{lemma}{引理}}{3.3}
a）设
\[0\longrightarrow F\longrightarrow G\longrightarrow H\]
是 $\mathcal C_S$ 中的\Term{exact-sequence}{正合列}。若 $G$ 为\Term{finite-type}{有限型}且 $H$ 为\Term{precoherent}{预凝聚}的，则 $F$ 为\Term{finite-type}{有限型}。
\SourcePage{123}
b）设
\[F\longrightarrow G\longrightarrow H\longrightarrow0\]
是 $\mathcal C_S$ 中的\Term{exact-sequence}{正合列}。若 $F$ 为\Term{finite-type}{有限型}且 $G$ 为\Term{precoherent}{预凝聚}的，则 $H$ 为\Term{precoherent}{预凝聚}的。

c）设
\[0\longrightarrow F\longrightarrow G\longrightarrow H\longrightarrow0\]
是 $\mathcal C_S$ 中的\Term{exact-sequence}{正合列}。若 $F$ 与 $H$ 都是\Term{precoherent}{预凝聚}的（相应地，\Term{finite-type}{有限型}的），则 $G$ 也如此。
\end{SGAresult}
\begin{proof}
a）作\Term{localization}{局部化}后，可取\Term{epimorphism}{满态射} $G'\to G$，其中 $G'\in\operatorname{ob}\mathcal C_0$。由此得到行\Term{exact}{正合}的\Term{commutative-diagram}{交换图}
\[\begin{tikzcd}
0\arrow[r]&F'\arrow[r]\arrow[d]&G'\arrow[r]\arrow[d]&H\arrow[d,"\mathrm{Id}"]\\
0\arrow[r]&F\arrow[r]&G\arrow[r]&H.
\end{tikzcd}\]
因为 $H$ 为\Term{precoherent}{预凝聚}的，$F'$ 为\Term{finite-type}{有限型}；而 $F'\to F$ 是\Term{epimorphism}{满态射}，所以 $F$ 为\Term{finite-type}{有限型}。

b）设 $U$ 在 $S$ 上方，且 $f:H'\to H|U$ 是 $\mathcal C_U$ 中的\Term{morphism}{态射}，其中 $H'\in\operatorname{ob}\mathcal C_{0U}$。改换记号后可设 $U=S$。由 $\mathcal C_0$ 在 $\mathcal C$ 中的\Term{local-scheme}{局部}\Term{quasi-liftable}{拟可提升性}，作\Term{localization}{局部化}后存在\Term{commutative-diagram}{交换方块}
\[\begin{tikzcd}G'\arrow[r,"p"]\arrow[d]&H'\arrow[d,"f"]\\G\arrow[r]&H,\end{tikzcd}\]
\SourcePage{124}
其中水平\Term{arrow}{箭头}都是\Term{epimorphism}{满态射}，且 $G'\in\operatorname{ob}\mathcal C_0$。因为 $\operatorname{Ker}(f)$ 是 $\operatorname{Ker}(fp)$ 的商，只须证明 $\operatorname{Ker}(fp)$ 为\Term{finite-type}{有限型}，所以可设 $H'=G'$。\Term{arrow}{箭头} $F\to G$ 与 $H'\to G$ 定义 $g:F\oplus H'\to G$。令 $K=\operatorname{Ker}(G\to H)$，便有下列\Term{commutative-diagram}{交换图}，其中各行\Term{exact}{正合}，$q$ 是\Term{epimorphism}{满态射}：
\[\begin{tikzcd}
0\arrow[r]&F\arrow[r]\arrow[d,"q"]&F\oplus H'\arrow[r]\arrow[d,"g"]&H'\arrow[r]\arrow[d,"f"]&0\\
0\arrow[r]&K\arrow[r]&G\arrow[r]&H\arrow[r]&0.
\end{tikzcd}\]
\Term{snake-lemma}{蛇引理}给出\Term{epimorphism}{满态射} $\operatorname{Ker}(g)\to\operatorname{Ker}(f)$。而 $F\oplus H'$ 为\Term{finite-type}{有限型}，$G$ 为\Term{precoherent}{预凝聚}的，故由 a），$\operatorname{Ker}(g)$ 为\Term{finite-type}{有限型}，结论即得。

c）先证明不带\Term{lie-bracket}{括号}的断言。改换记号后，所需证明的是：若 $g:G'\to G$ 是 $\mathcal C_S$ 中的\Term{morphism}{态射}，$G'\in\operatorname{ob}\mathcal C_0$，则 $\operatorname{Ker}(g)$ 为\Term{finite-type}{有限型}。以 $H'$ 表示复合\Term{arrow}{箭头} $G'\to H$ 的像，有行\Term{exact}{正合}的\Term{commutative-diagram}{交换图}
\[\begin{tikzcd}
0\arrow[r]&F'\arrow[r]\arrow[d,"f"]&G'\arrow[r]\arrow[d,"g"]&H'\arrow[r]\arrow[d,"h"]&0\\
0\arrow[r]&F\arrow[r]&G\arrow[r]&H\arrow[r]&0.
\end{tikzcd}\]
由于 $h$ 为\Term{monomorphism}{单态射}，\Term{snake-lemma}{蛇引理}给出\Term{isomorphism}{同构} $\operatorname{Ker}(f)\simeq\operatorname{Ker}(g)$。但由 a），$F'$ 为\Term{finite-type}{有限型}；再次应用 a）可知 $\operatorname{Ker}(f)$ 为\Term{finite-type}{有限型}，从而得到不带\Term{lie-bracket}{括号}的断言。再证\Term{lie-bracket}{括号}中的断言（严格说来，此处不能引用（2.5 b）），因为没有假定 $\mathcal C_0$ \Term{local-scheme}{局部}地对\Term{epimorphism}{满态射}的核\Term{quasi-stable}{拟稳定}）。\Term{localization}{局部化}后存在\Term{epimorphism}{满态射} $u:F'\to F$、$v:H'\to H$，其中
\SourcePage{125}
$F',H'\in\operatorname{ob}\mathcal C_0$。进一步\Term{localization}{局部化}后可设 $v$ 经过 $G$ \Term{resolution}{分解}，因而有行\Term{exact}{正合}的\Term{commutative-diagram}{交换图}
\[\begin{tikzcd}
0\arrow[r]&F'\arrow[r]\arrow[d,"u"]&F'\oplus H'\arrow[r]\arrow[d,"w"]&H'\arrow[r]\arrow[d,"v"]&0\\
0\arrow[r]&F\arrow[r]&G\arrow[r]&H\arrow[r]&0.
\end{tikzcd}\]
因为 $u,v$ 都是\Term{epimorphism}{满态射}，$w$ 也是，结论即得。这就完成了（3.3）的证明。
\end{proof}

立即得到下列推论。
\begin{SGAresult}{推论}{3.4}
$\mathcal C_S$ 中由\Term{coherent}{凝聚}\Term{object}{对象}组成的（严格）\Term{full-subcategory}{满子范畴}对\Term{finite}{有限}\Term{projective-limit}{射影极限}（相应地，\Term{inductive-limit}{归纳极限}）\Term{stable}{稳定}。特别地，两个\Term{coherent}{凝聚}\Term{object}{对象}之间\Term{morphism}{态射}的核、\Term{cokernel}{余核}和像都是\Term{coherent}{凝聚}的。设
\[E^0\longrightarrow E^1\longrightarrow E^2\longrightarrow E^3\longrightarrow E^4\]
是 $\mathcal C_S$ 中的\Term{exact-sequence}{正合列}。若 $E^0,E^1,E^3,E^4$ 都\Term{coherent}{凝聚}，则 $E^2$ 也\Term{coherent}{凝聚}。
\end{SGAresult}

\begin{SGAresult}{推论}{3.5}
假定 $\mathcal C_0$ 的\Term{object}{对象}都是 $\mathcal C_0$-\Term{coherent}{凝聚}的（这特别蕴含 $\mathcal C_0$ \Term{local-scheme}{局部}地对\Term{epimorphism}{满态射}的核\Term{quasi-stable}{拟稳定}，即满足（2.0）的条件）。

a）对于 $F\in\operatorname{ob}\mathcal C_S$，下列条件\Term{equivalence}{等价}：
\begin{enumerate}
\item[(i)] $F$ 为\Term{finite-presentation}{有限表示}的；
\item[(ii)] $F$ 为\Term{coherent}{凝聚}的；
\item[(iii)] $F$ 为 $\infty$-\Term{finite-presentation}{有限表示}的（2.8）。
\end{enumerate}
b）设 $F\in\operatorname{ob}D^-(\mathcal C_S)$，$n\in\mathbb Z$。则
$F$ 为 $n$-\Term{pseudo-coherent}{伪凝聚}的，当且仅当 $i\geq n+1$ 时 $H^i(F)$ 都\Term{coherent}{凝聚}，且 $H^n(F)$ 为\Term{finite-type}{有限型}。
\SourcePage{126}
特别地，$F$ 为\Term{pseudo-coherent}{伪凝聚}的，当且仅当每个 $H^n(F)$ 都\Term{coherent}{凝聚}。
\end{SGAresult}
\begin{proof}
a）以及 b）的正向蕴含显然由（3.4）得到。b）的反向蕴含立即由 a）及（2.11 b））得到。
\end{proof}

\textbf{例 3.6。}设 $(\mathscr S,A)$ 为\Term{ringed-topos}{赋环拓扑斯}，$\mathcal C,\mathcal C_0$ 为（1.3 d））中的 $\mathscr S$-\Term{category}{范畴}（亦见（2.15））。$\mathcal C_0$ 的\Term{object}{对象}都\Term{coherent}{凝聚}，当且仅当 $A$ \Term{coherent}{凝聚}（即\Term{precoherent}{预凝聚}）。以下是满足此条件的例子：
\begin{itemize}
\item $A$ 是值为左 Noether 环的\Term{constant-sheaf}{常值层}（SGA 4 IX 2.1）。
\item $\mathscr S$ 是\Term{locally-noetherian}{局部 Noether} \Term{scheme}{概形} $S$ 上（\Term{zariski-topology}{Zariski 拓扑}下的）\Term{set}{集合}\Term{sheaf}{层}构成的\Term{topos}{拓扑斯}\footnote{依照新的约定，我们将用“\Term{scheme}{概形}”代替“\Term{prescheme}{预概形}”，用“\Term{separated-scheme}{分离概形}”代替原先的“\Term{scheme}{概形}”。}，$A=\mathcal O_S$。
\item $\mathscr S$ 是\Term{complex-analytic-space}{复解析空间} $S$ 上的\Term{set}{集合}\Term{sheaf}{层}构成的\Term{topos}{拓扑斯}，$A=\mathcal O_S$。
\end{itemize}
另一方面，\Term{topological-space}{拓扑空间}上的\Term{continuous}{连续}复值函数\Term{sheaf}{层}\Term{generic}{一般}并不\Term{coherent}{凝聚}。

当然，若 $(\mathscr S,A)$ 由一个\Term{ringed-space}{赋环空间}定义，则这里的\Term{coherent}{凝聚性}与 Serre 的通常概念一致（EGA $0_{\mathrm I}$ 5.3.1）。

\begin{SGAresult}{命题}{3.7}
a）设 $f:(\mathscr S',A')\to(\mathscr S,A)$ 是\Term{ringed-topos}{赋环拓扑斯}的\Term{morphism}{态射}。假定 $f$ 右\Term{flat}{平坦}，即\Term{functor}{函子} $f^*=A'\otimes_A f^{-1}(\,\cdot\,)$ \Term{exact}{正合}。对于 $F\in\operatorname{ob}D({}_A\mathscr S)$、$G\in\operatorname{ob}D^+({}_A\mathscr S)$，有\Term{canonical}{典范}的、\Term{functorial}{函子性}的\Term{arrow}{箭头}
\[
f^*\operatorname{R\mathcal Hom}(F,G)\longrightarrow
\operatorname{R\mathcal Hom}(f^*F,f^*G).
\]
\SourcePage{127}
当 $F$ 为\Term{pseudo-coherent}{伪凝聚}的时，此\Term{arrow}{箭头}是\Term{isomorphism}{同构}。

b）假定 $A$ 交换且\Term{coherent}{凝聚}，设 $F\in\operatorname{ob}D^-(\mathscr S)$、$G\in\operatorname{ob}D^+(\mathscr S)$。若每个 $H^i(F),H^i(G)$ 都\Term{coherent}{凝聚}，则每个 $\mathcal Ext^i(F,G)$ 也\Term{coherent}{凝聚}。
\end{SGAresult}
\begin{proof}
a）为定义此\Term{arrow}{箭头}（参看 SGA 4 XVII），可设 $G$ 各项均为\Term{injective}{内射}\Term{object}{对象}，再取复合\Term{arrow}{箭头}
\[
f^*\mathcal Hom(F,G)\longrightarrow\mathcal Hom(f^*F,f^*G)
\longrightarrow\operatorname{R\mathcal Hom}(f^*F,f^*G).
\]
证明当 $F$ 为\Term{pseudo-coherent}{伪凝聚}的时它是\Term{isomorphism}{同构}。为简便，以 $T$（相应地，$T'$）表示\Term{functor}{函子} $f^*\operatorname{R\mathcal Hom}(\,\cdot\,,G)$（相应地，$\operatorname{R\mathcal Hom}(f^*(\,\cdot\,),f^*G)$）。取 $a\in\mathbb Z$，使 $i<a$ 时 $H^i(G)=0$。显然 $TA\simeq T'A$，故通过逐步\Term{resolution}{分解}，对于严格\Term{perfect-complex}{完美}的 $F$ 有 $TF\simeq T'F$（2.1）。另一方面，若 $F$ 在\Term{degree}{次数} $>n$ \Term{acyclic}{无上同调}，则 $TF,T'F$ 在\Term{degree}{次数} $<a-n$ \Term{acyclic}{无上同调}。为验证当 $F$ 为\Term{pseudo-coherent}{伪凝聚}的时 $TF\simeq T'F$，由于问题是\Term{local-scheme}{局部}的，可设 $F$ 为严格 $n$-\Term{pseudo-coherent}{伪凝聚}，即有\Term{distinguished-triangle}{可区别三角形}
\[F'\longrightarrow F\longrightarrow F''\xrightarrow{+1}F',\]
其中 $F'$ 严格\Term{perfect-complex}{完美}，$F''$ 在\Term{degree}{次数} $>n-1$ \Term{acyclic}{无上同调}。由此得到\Term{distinguished-triangle}{可区别三角形}的\Term{morphism}{态射}
\[\begin{tikzcd}
TF''\arrow[r]\arrow[d]&TF\arrow[r]\arrow[d]&TF'\arrow[r,"{-1}"]\arrow[d]&TF''[-1]\arrow[d]\\
T'F''\arrow[r]&T'F\arrow[r]&T'F'\arrow[r,"{-1}"]&T'F''[-1].
\end{tikzcd}\]
由于 $F'$ 严格\Term{perfect-complex}{完美}，$TF'\simeq T'F'$；而 $TF''$ 和 $T'F''$
\SourcePage{128}
在\Term{degree}{次数} $\leq a-n$ \Term{acyclic}{无上同调}。写出上述三角形的\Term{cohomology}{上同调}\Term{exact-sequence}{正合列}之间的诱\Term{derived-morphism}{导态射}，可知 $TF\to T'F$ 在 $i\leq a-n$ 时诱导\Term{isomorphism}{同构} $H^i(TF)\simeq H^i(T'F)$。由于 $n$ 可以取得任意负，结论即得。

b）证明留给读者作练习；例如参看（[H] II 3.3）。
\end{proof}

\textbf{注记 3.7.1。}稍后（7.1.2）还将看到 $\operatorname{R\mathcal Hom}$ 与 $\mathrm Lf^*$ “相容”的另一个情形。

\begin{SGAresult}{推论}{3.7.2}
（参看 EGA $0_{\mathrm I}$ 5.2.7。）设 $(\mathscr S,A)$ 为\Term{ringed-topos}{赋环拓扑斯}，$s$ 为 $\mathscr S$ 的点（SGA 4 VIII 7.8），$E,F\in\operatorname{ob}D^b({}_A\mathscr S)_{\mathrm{coh}}$。若 $E_s\simeq F_s$，则存在 $\mathscr S$ 中的 $s$-\Term{pointed-object}{带点对象} $U$，使 $E|U\simeq F|U$。
\end{SGAresult}
\begin{proof}
设 $\alpha:E_s\to F_s$ 与 $\beta:F_s\to E_s$ 是互逆的\Term{isomorphism}{同构}。由（2.22 a）），有
\[
\mathcal Hom(E,F)_s\simeq\operatorname{Hom}(E_s,F_s)
\quad\text{（相应地，}\mathcal Hom(F,E)_s\simeq\operatorname{Hom}(F_s,E_s)\text{）。}
\]
故存在 $s$ 的一个“\Term{neighborhood}{邻域}” $U$ 及 $U$ 上 $\mathcal Hom(E,F)$（相应地，$\mathcal Hom(F,E)$）的\Term{section}{截面} $u$（相应地，$v$），使 $u_s=\alpha$（相应地，$v_s=\beta$）。\Term{homomorphism}{同态} $vu:E|U\to E|U$ 与 $uv:F|U\to F|U$ 在 $s$ 处诱导\Term{identity-morphism}{恒等态射}，因此在 $U$ 上方某个适当的 $V$ 上也诱导\Term{identity-morphism}{恒等态射}，证毕。
\end{proof}

\SourcePage{129}
\SGAsection{4}{完美复形}
\textbf{4.0。}本节中，$\mathscr S$ 表示一个\Term{site}{景}，$S$ 为 $\mathscr S$ 的\Term{object}{对象}，$\mathcal C$ 为\Term{flat}{平坦}的 Abel $\mathscr S$-\Term{category}{范畴}，$\mathcal C_0$ 为 $\mathcal C$ 的严格\Term{full-subcategory}{满}\Term{additive}{加性}子 $\mathscr S$-\Term{category}{范畴}（1.1）。除非另有说明，假定 $\mathcal C_0$ 在 $\mathcal C$ 中\Term{local-scheme}{局部}\Term{liftable}{可提升}，并\Term{local-scheme}{局部}地对\Term{epimorphism}{满态射}的核\Term{stable}{稳定}（1.2）。

\begin{SGAresult}{\Term{lemma}{引理}}{4.1}
\footnote{此\Term{lemma}{引理}由 D. Ferrand 告知编写者。}设 $u:E\to F$ 是 $C(\mathcal C_S)$ 中的\Term{arrow}{箭头}。若 $E$ 为严格 $\mathcal C_0$-\Term{perfect-complex}{完美}的（2.1），且 $F$ \Term{acyclic}{无上同调}，则 $u$ \Term{local-scheme}{局部}地\Term{homotopy}{同伦}于零。
\end{SGAresult}
\begin{proof}
需\Term{local-scheme}{局部}构造\Term{homotopy}{同伦} $k$，使 $u=dk+kd$。因为 $E$ 严格\Term{perfect-complex}{完美}，存在 $a,b\in\mathbb Z$，$a\leq b$，使 $i\notin[a,b]$ 时 $E^i=0$。因而 $i\notin[a,b]$ 时可取 $k^i=0$。随后对 $i$ 作降归纳，为 $i\geq a$ 构造 $k^i$。假设已构造 $k^i:E^i\to F^{i-1}$，使得 $i\geq n>a$ 时
\[u^i=k^{i+1}d_E^i+d_F^{i-1}k^i.\]
证明作\Term{localization}{局部化}后存在 $k^{n-1}:E^{n-1}\to F^{n-2}$，使
\[u^{n-1}=k^nd_E^{n-1}+d_F^{n-2}k^{n-1}.\]
由归纳假设，\Term{arrow}{箭头}
\[v^{n-1}=u^{n-1}-k^nd_E^{n-1}:E^{n-1}\longrightarrow F^{n-1}\]
经过 $Z^{n-1}(F)$ \Term{resolution}{分解}。而 $F$ \Term{acyclic}{无上同调}，所以 $Z^{n-1}(F)=B^{n-1}(F)$。由 $\mathcal C_0$ 在 $\mathcal C$ 中\Term{local-scheme}{局部}\Term{liftable}{可提升}，作\Term{localization}{局部化}后存在 $k^{n-1}:E^{n-1}\to F^{n-2}$，使 $v^{n-1}=d_F^{n-2}k^{n-1}$，证毕。
\end{proof}

\begin{SGAresult}{推论}{4.2}
设
\[\begin{tikzcd}&P\arrow[d,"v"]\\E\arrow[r,"f"]&F\end{tikzcd}\]
\SourcePage{130}
为 $C(\mathcal C_S)$ 中的图，其中 $P$ 为严格 $\mathcal C_0$-\Term{perfect-complex}{完美}的，$f$ 是\Term{quasi-isomorphism}{拟同构}。$S$ 上方满足下列条件的 $X$ 的\Term{set}{集合}是 $S$ 的一个\Term{refinement}{加细}：存在 $C(\mathcal C_X)$ 中的\Term{arrow}{箭头} $u:P_X\to E_X$，使下图在 $K(\mathcal C_X)$ 中交换：
\[\begin{tikzcd}&P_X\arrow[dl,"u"']\arrow[d,"v_X"]\\E_X\arrow[r,"f_X"]&F_X.\end{tikzcd}\]
\end{SGAresult}
\begin{proof}
设 $G$ 为 $f$ 的\Term{cone}{锥}。有\Term{exact-sequence}{正合列}
\[\cdots\longrightarrow\operatorname{Hom}_{K(\mathcal C_S)}(P,E)
\longrightarrow\operatorname{Hom}_{K(\mathcal C_S)}(P,F)
\longrightarrow\operatorname{Hom}_{K(\mathcal C_S)}(P,G)\longrightarrow\cdots.\]
因为 $G$ \Term{acyclic}{无上同调}，由（4.1），$v$ 在 $\operatorname{Hom}_{K(\mathcal C_S)}(P,G)$ 中的像\Term{local-scheme}{局部}为零。故作\Term{localization}{局部化}后，存在 $u:P\to E$，使 $v=fu$ 在 $K(\mathcal C_S)$ 中成立，证毕。
\end{proof}

\begin{SGAresult}{推论}{4.3}
设 $f:E\to F$ 是 $C(\mathcal C_S)$ 中的\Term{quasi-isomorphism}{拟同构}，$F$ 严格\Term{perfect-complex}{完美}（相对于 $\mathcal C_0$）。$S$ 上方满足下列条件的 $X$ 的\Term{set}{集合}是 $S$ 的一个\Term{refinement}{加细}：存在\Term{quasi-isomorphism}{拟同构} $g:F_X\to E_X$，使 $f_Xg$ \Term{homotopy}{同伦}于 $\mathrm{Id}_{F_X}$。
\end{SGAresult}
\begin{proof}这是（4.2）的一个特例。\end{proof}

\textbf{评注 4.4。}设 $F\in\operatorname{ob}K(\mathcal C_S)$。可以定义一个在 $\mathscr S_{/S}$ 上的\Term{fibered-category}{纤维范畴} $(\mathrm{Qis}/F)$（例如借助 $\mathcal C$ 的一个\Term{cleavage}{劈裂}，见 SGA 1 VI 7.1），它在 $S$ 上方的 $X$ 处的\Term{fiber}{纤维}是\Term{category}{范畴} $(\mathrm{Qis}/F_X)^\circ$（以 $F_X$ 为目标的\Term{quasi-isomorphism}{拟同构}所成\Term{category}{范畴}的\Term{opposite-category}{反范畴}）。于是（4.3）也可表述为：
\SourcePage{131}
若 $F$ 严格\Term{perfect-complex}{完美}，则\Term{arrow}{箭头} $\mathrm{Id}_F$ 在 $(\mathrm{Qis}/F)$ 中\Term{local-scheme}{局部}\Term{cofinal}{共尾}。

\textbf{记号 4.5。}设 $\mathcal M$ 为 $\mathscr S$-\Term{fibered-category}{纤维范畴}，$X\in\operatorname{ob}\mathscr S$，$E,F\in\operatorname{ob}\mathcal M_X$。以 $\mathcal Hom_{\mathcal M}(E,F)$ 表示 $\mathscr S_{/X}$ 上与 $\mathcal M$ 中从 $E$ 到 $F$ 的 $\mathscr S$-\Term{morphism}{态射}\Term{presheaf}{预层}相伴的\Term{sheaf}{层}（此\Term{presheaf}{预层}的定义见 [G] I 2.6；应当注意，即使 $\mathcal C$ 是\Term{stack}{叠}（[G]），取 $\mathcal M=K(\mathcal C)$ 或 $D(\mathcal C)$ 时所对应的\Term{presheaf}{预层}\Term{generic}{一般}也不是\Term{sheaf}{层}）。

\begin{SGAresult}{推论}{4.6}
设 $E,F\in\operatorname{ob}C(\mathcal C_S)$，其中 $E$ 严格\Term{perfect-complex}{完美}。\Term{canonical}{典范}\Term{arrow}{箭头}
\[\mathcal Hom_{K(\mathcal C)}(E,F)\longrightarrow\mathcal Hom_{D(\mathcal C)}(E,F)\]
是\Term{isomorphism}{同构}。
\end{SGAresult}
\begin{proof}
这立即由（4.4）及对 $S$ 上方的 $X$ 成立的公式
\[
\operatorname{Hom}_{D(\mathcal C_X)}(E_X,F_X)
=\varinjlim_{(\mathrm{Qis}/E_X)^\circ}\operatorname{Hom}_{K(\mathcal C_X)}(\,\cdot\,,F_X)
\]
得到。
\end{proof}

\begin{SGAresult}{定义}{4.7}
设 $F\in\operatorname{ob}D(\mathcal C_S)$，$a,b\in\mathbb Z$，$a\leq b$。若 $S$ 上方满足下列条件的 $X$ 的\Term{set}{集合}是 $S$ 的一个\Term{refinement}{加细}：存在
\SourcePage{132}
$C(\mathcal C_X)$ 中的\Term{quasi-isomorphism}{拟同构} $F'\to F$，其中 $F'$ 严格 $\mathcal C_0$-\Term{perfect-complex}{完美}（2.1），且 $i\notin[a,b]$ 时 $F'^i=0$，则称 $F$ 的 $\mathcal C_0$-\Term{perfect-amplitude}{完美振幅}包含于区间 $[a,b]$，并写作
\[\mathcal C_0\text{-}\operatorname{parf\!\! -amp}(F)\subset[a,b].\]
（由（4.3），将“$C(\mathcal C_X)$ 中的\Term{quasi-isomorphism}{拟同构}”改为“$D(\mathcal C_X)$ 中的\Term{isomorphism}{同构}”可得\Term{equivalence}{等价}条件。）若存在整数 $a,b$ 使 $\mathcal C_0\text{-}\operatorname{parf\!\! -amp}(F)\subset[a,b]$，则称 $F$ 的 $\mathcal C_0$-\Term{perfect-amplitude}{完美振幅}\Term{finite}{有限}。若 $F$ \Term{local-scheme}{局部}具有\Term{finite}{有限}的 $\mathcal C_0$-\Term{perfect-amplitude}{完美振幅}，则称 $F$ 是 $\mathcal C_0$-\Term{perfect-complex}{完美}的。\footnote{当然，在不会引起混淆时，上述表达式中的 $\mathcal C_0$ 将予省略。}
\end{SGAresult}

\textbf{例 4.8。}设 $(\mathscr S,A)$ 为\Term{ringed-topos}{赋环拓扑斯}；$\mathcal C$ 为 $U\mapsto$（左 $A_U$-模的\Term{category}{范畴}）定义的 $\mathscr S$-\Term{category}{范畴}；$\mathcal C_0$ 为 $U\mapsto$（作为\Term{finite-free}{有限型自由模}之\Term{direct-summand}{直和因子}的左 $A_U$-模的\Term{category}{范畴}）定义的子 $\mathscr S$-\Term{category}{范畴}。由（1.3.1），$(\mathcal C,\mathcal C_0)$ 满足（4.0）的假设。若\Term{object}{对象} $F\in D({}_A\mathscr S)$（参看（2.15））的 $\mathcal C_0$-\Term{perfect-amplitude}{完美振幅}包含于 $[a,b]$（相应地，$\mathcal C_0$-\Term{perfect-amplitude}{完美振幅}\Term{finite}{有限}，或 $F$ 为 $\mathcal C_0$-\Term{perfect-complex}{完美}的），则称 $F$ 的\Term{perfect-amplitude}{完美振幅}包含于 $[a,b]$（相应地，\Term{perfect-amplitude}{完美振幅}\Term{finite}{有限}，或 $F$ 为\Term{perfect-complex}{完美}的）。当\Term{finite-free}{有限型自由模}的\Term{direct-summand}{直和因子}\Term{locally-free}{局部自由}时（参看（2.15.1）），说 $F$ \Term{perfect-complex}{完美}，即是说 $F$ 在 $D(\mathscr S)$ 中\Term{local-scheme}{局部}\Term{isomorphism}{同构}于\Term{degree}{次数}\Term{bounded}{有界}、各项为\Term{finite-free}{有限型自由模}的\Term{complex}{复形}。

\textbf{记号 4.9。}用 $D(\mathcal C)_{\mathcal C_0\text{-parf}}$ 或 $\operatorname{Parf}(\mathcal C,\mathcal C_0)$ 表示 $D(\mathcal C)$ 中由 $\mathcal C_0$-\Term{perfect-complex}{完美}\Term{object}{对象}构成的严格\Term{full-subcategory}{满}子 $\mathscr S$-\Term{category}{范畴}。在例（4.8）的情形，记作 $D({}_A\mathscr S)_{\mathrm{parf}}$、$\operatorname{Parf}({}_A\mathscr S)$，或在 $A$ 不会引起混淆时记作 $\operatorname{Parf}(\mathscr S)$。

\SourcePage{133}
\begin{SGAresult}{命题}{4.10}
\Term{category}{范畴} $\operatorname{Parf}(\mathcal C,\mathcal C_0)$ 是 $D(\mathcal C)$ 的\Term{triangulated-subcategory}{三角}子 $\mathscr S$-\Term{category}{范畴}（1.1.4）。\Term{canonical}{典范}\Term{functor}{函子}
\[K^b(\mathcal C_0)\longrightarrow\operatorname{Parf}(\mathcal C,\mathcal C_0)\]
（相应地，$\operatorname{Tr}(K^b(\mathcal C_0))\to\operatorname{Tr}(\operatorname{Parf}(\mathcal C,\mathcal C_0))$，见（1.1.3））在相伴的\Term{stack}{叠}上诱导 $\mathscr S$-\Term{equivalence-of-categories}{等价}（[G] II 2.2.4），亦即：

（i）对每个 $X\in\operatorname{ob}\mathscr S$ 及所有 $E,F\in\operatorname{ob}K^b(\mathcal C_{0X})$（相应地，$\operatorname{Tr}(K^b(\mathcal C_{0X}))$），\Term{canonical}{典范}\Term{arrow}{箭头}
\[\mathcal Hom_{K^b(\mathcal C_0)}(E,F)\longrightarrow\mathcal Hom_{D(\mathcal C)}(E,F)\]
（相应地，$\mathcal Hom_{\operatorname{Tr}(K^b(\mathcal C_0))}(E,F)\to\mathcal Hom_{\operatorname{Tr}(D(\mathcal C))}(E,F)$）是\Term{isomorphism}{同构}；

（ii）$\operatorname{Parf}(\mathcal C,\mathcal C_0)$（相应地，$\operatorname{Tr}(\operatorname{Parf}(\mathcal C,\mathcal C_0))$）的每个\Term{object}{对象}，在 $D(\mathcal C)$（相应地，$\operatorname{Tr}(D(\mathcal C))$）中\Term{local-scheme}{局部}\Term{isomorphism}{同构}于 $K^b(\mathcal C_0)$（相应地，$\operatorname{Tr}(K^b(\mathcal C_0))$）的\Term{object}{对象}。
\end{SGAresult}
\begin{proof}
（i）不带\Term{lie-bracket}{括号}的断言是（4.6）的特例；（ii）不带\Term{lie-bracket}{括号}的断言由定义成立（（4.7）和（4.9））。另一方面，$\operatorname{Parf}(\mathcal C,\mathcal C_0)$ 显然对\Term{degree}{次数}\Term{translation}{平移}\Term{stable}{稳定}。因此，要证明它是 $D(\mathcal C)$ 的\Term{triangulated-subcategory}{三角}子 $\mathscr S$-\Term{category}{范畴}，只须证明：若
\[E\xrightarrow{u}F\longrightarrow G\xrightarrow{+1}E\tag{*}\]
是 $D(\mathcal C)$ 中的\Term{distinguished-triangle}{可区别三角形}，且 $E,F$ \Term{perfect-complex}{完美}，则 $G$
\SourcePage{134}
也\Term{perfect-complex}{完美}。\Term{localization}{局部化}后，可设 $E,F$ 为 $K^b(\mathcal C_0)$ 的\Term{object}{对象}。由（i）不带\Term{lie-bracket}{括号}的断言，进一步\Term{localization}{局部化}后可设 $u$ 为 $K^b(\mathcal C_0)$ 中的\Term{arrow}{箭头}；这时 $G$ \Term{isomorphism}{同构}于 $u$ 的\Term{cone}{锥}，后者严格\Term{perfect-complex}{完美}，所以 $G$ \Term{perfect-complex}{完美}。故 $\operatorname{Parf}(\mathcal C,\mathcal C_0)$ 确为 $D(\mathcal C)$ 的\Term{triangulated-subcategory}{三角}子 $\mathscr S$-\Term{category}{范畴}。上述论证同时证明了（ii）\Term{lie-bracket}{括号}中的断言。

只须再证明（i）\Term{lie-bracket}{括号}中的断言。设 $X\in\operatorname{ob}\mathscr S$，
\[
E=(E_1\xrightarrow{a_1}E_2\xrightarrow{a_2}E_3\xrightarrow[+1]{a_3}E_1),\qquad
F=(F_1\xrightarrow{b_1}F_2\xrightarrow{b_2}F_3\xrightarrow[+1]{b_3}F_1)
\]
为 $K^b(\mathcal C_{0X})$ 的\Term{distinguished-triangle}{可区别三角形}。需证明\Term{arrow}{箭头}
\[\mathcal Hom_{\operatorname{Tr}(K^b(\mathcal C_0))}(E,F)\longrightarrow
\mathcal Hom_{\operatorname{Tr}(D(\mathcal C))}(E,F)\tag{**}\]
是\Term{isomorphism}{同构}。由（i）不带\Term{lie-bracket}{括号}的断言，已知（**）为\Term{injectivity}{单射}。设 $(u_i):E\to F$ 为 $D(\mathcal C)$ 中三角形的\Term{morphism}{态射}。由（i）不带\Term{lie-bracket}{括号}的断言，\Term{localization}{局部化}后可设各 $u_i$（$1\leq i\leq3$）为 $K^b(\mathcal C_0)$ 中的\Term{arrow}{箭头}。对 $1\leq i\leq3$，\Term{arrow}{箭头} $u_{i+1}a_i-b_iu_i$（约定 $u_4=u_1$）在 $D(\mathcal C)$ 中为零，故仍由（i）不带\Term{lie-bracket}{括号}的断言，它在 $K^b(\mathcal C_0)$ 中\Term{local-scheme}{局部}为零。换言之，$(u_i)$ \Term{local-scheme}{局部}定义 $K^b(\mathcal C_0)$ 中三角形的\Term{morphism}{态射}。这就证明了（**）为\Term{surjectivity}{满射}，并完成（4.10）的证明。
\end{proof}

\textbf{注记 4.10.1。}在（4.10）中，可以把 $\operatorname{Parf}(\mathcal C,\mathcal C_0)$ 换成
\[D^*(\mathcal C)_{\mathcal C_0\text{-parf}}=D^*(\mathcal C)\cap\operatorname{Parf}(\mathcal C,\mathcal C_0)\qquad(*=-\text{ 或 }b),\]
证明无须更改。

\textbf{补充 4.11。}可以用\Term{perfect-amplitude}{完美振幅}（4.7）更精确地描述 $\operatorname{Parf}(\mathcal C,\mathcal C_0)$ 的三角结构。若 $E\in\operatorname{ob}D(\mathcal C)$ 的\Term{perfect-amplitude}{完美振幅}包含于 $[a,b]$，则 $E[n]$ 的\Term{perfect-amplitude}{完美振幅}包含于 $[a-n,b-n]$。
\SourcePage{135}
设
\[E\longrightarrow F\longrightarrow G\xrightarrow{+1}E\]
为 $D(\mathcal C)$ 的\Term{distinguished-triangle}{可区别三角形}。可列出与（2.6）类似的表格，表中各格给出\Term{perfect-amplitude}{完美振幅}，每一行对应一个蕴含，结论以下划线标出。
\[
\begin{array}{ccc}
E&F&G\\[2pt]
[a+1,b+1]&[a,b]&\underline{[a,b]}\\
{}[a,b]&\underline{[a,b]}&[a,b]\\
\underline{[a,b]}&[a,b]&[a-1,b-1]
\end{array}
\]
由（4.10）立即得到下列推论。

\begin{SGAresult}{推论}{4.12}
假定 $\mathscr S$ 为\Term{point-site}{点景}。$F\in\operatorname{ob}D(\mathcal C)$ 的\Term{perfect-amplitude}{完美振幅}（相对于 $\mathcal C_0$）包含于 $[a,b]$，当且仅当存在\Term{quasi-isomorphism}{拟同构} $F'\to F$，其中 $F'$ 严格\Term{perfect-complex}{完美}，且 $i\notin[a,b]$ 时 $F'^i=0$。\Term{canonical}{典范}\Term{functor}{函子}
\[K^b(\mathcal C_0)\longrightarrow\operatorname{Parf}(\mathcal C,\mathcal C_0)\]
及
\[\operatorname{Tr}(K^b(\mathcal C_0))\longrightarrow\operatorname{Tr}(\operatorname{Parf}(\mathcal C,\mathcal C_0))\]
都是\Term{equivalence-of-categories}{范畴等价}。
\end{SGAresult}

\textbf{注记 4.12.1。}当 $\mathscr S$ 为\Term{point-site}{点景}时，$\mathcal C_0$ 的\Term{object}{对象}都是\Term{projective-object}{投射}的。当 $\mathcal C_0$ 包含 $\mathcal C$ 的全部\Term{projective-object}{投射}\Term{object}{对象}时，将“$\mathcal C_0$-\Term{perfect-amplitude}{完美振幅}”称为“\Term{projective-amplitude}{投射振幅}”。

\SourcePage{136}
\textbf{注记 4.12.2。}设 $\mathcal C$ 为\Term{abelian-category}{Abel 范畴}（可视为\Term{point-site}{点景}上的\Term{fibered-category}{纤维范畴}），$\mathcal C_0$ 为其\Term{additive}{加性}严格\Term{full-subcategory}{满子范畴}。假定 $\mathcal C_0$ 对\Term{epimorphism}{满态射}的核\Term{stable}{稳定}，且在 $\mathcal C$ 中\Term{quasi-liftable}{拟可提升}。这时 $K^b(\mathcal C_0)$ 中\Term{acyclic}{无上同调}的\Term{object}{对象}\Term{generic}{一般}并不为零，故不能期待\Term{functor}{函子} $K^b(\mathcal C_0)\to D(\mathcal C)$ 是\Term{faithful}{忠实}的。不过，以 $K^{b,\phi}(\mathcal C_0)$ 表示 $K^b(\mathcal C_0)$ 中由（在 $\mathcal C$ 中）\Term{acyclic}{无上同调}的\Term{object}{对象}构成的\Term{thick-subcategory}{厚子范畴}，则\Term{canonical}{典范}\Term{functor}{函子}
\[K^b(\mathcal C_0)/K^{b,\phi}(\mathcal C_0)\longrightarrow D(\mathcal C)\]
是\Term{fully-faithful}{全忠实}的。

提示：已经知道（2.7），\Term{functor}{函子}
\[K^-(\mathcal C_0)/K^{-,\phi}(\mathcal C_0)\longrightarrow D^-(\mathcal C)\]
是\Term{fully-faithful}{全忠实}的。证明（例如使用判据（[V] 4.2 b）（ii）$'$）\Term{functor}{函子}
\[K^b(\mathcal C_0)/K^{b,\phi}(\mathcal C_0)\longrightarrow K^-(\mathcal C_0)/K^{-,\phi}(\mathcal C_0)\]
是\Term{fully-faithful}{全忠实}的。

\begin{SGAresult}{\Term{lemma}{引理}}{4.13}
设 $F\in\operatorname{ob}D(\mathcal C_S)$，$a,b,c\in\mathbb Z$，$a\leq b$ 且 $a\leq c$。若 $F$ 的\Term{perfect-amplitude}{完美振幅}包含于 $[a,c]$（4.7），且 $i>b$ 时 $H^i(F)=0$，则 $F$ 的\Term{perfect-amplitude}{完美振幅}包含于 $[a,d]$，其中 $d=\inf(b,c)$。
\end{SGAresult}
\begin{proof}
可设 $b\leq c$，且 $F$ 严格\Term{perfect-complex}{完美}，$i\notin[a,c]$ 时 $F^i=0$。于是 $F$ \Term{isomorphism}{同构}于其\Term{truncation}{截断} $F'$，后者定义为
\SourcePage{137}
\[F'^i=F^i\ (i<b),\qquad F'^b=Z^b(F),\qquad F'^i=0\ (i>b).\]
而 $Z^b(F)$ 有一个由 $\mathcal C_0$ 的\Term{object}{对象}构成的\Term{finite}{有限}右\Term{resolution}{消解}，所以它\Term{local-scheme}{局部}为 $\mathcal C_0$ 的\Term{object}{对象}，断言即得。
\end{proof}

\begin{SGAresult}{推论}{4.14}
设 $F\in\operatorname{ob}\mathcal C_S$，$n\in\mathbb N$。下列条件\Term{equivalence}{等价}：

（i）$\operatorname{parf.amp}(F)\subset[-n,0]$；

（ii）$F$ \Term{local-scheme}{局部}有一个由 $\mathcal C_0$ 的\Term{object}{对象}构成、\Term{length}{长度}为 $n$ 的左\Term{resolution}{消解}，即 $S$ 上方满足下列条件的 $X$ 的\Term{set}{集合}是 $S$ 的一个\Term{refinement}{加细}：存在\Term{exact-sequence}{正合列}
\[0\longrightarrow L^{-n}\longrightarrow L^{-n+1}\longrightarrow\cdots\longrightarrow L^0\longrightarrow F_X\longrightarrow0,\]
其中对每个 $i$ 都有 $L^i\in\operatorname{ob}\mathcal C_{0X}$。
\end{SGAresult}
\begin{proof}这是（4.7）与（4.13）的直接推论。\end{proof}

\begin{SGAresult}{定义}{4.14.1}
若 $F$ 满足（4.14）的\Term{equivalence}{等价}条件，则称 $F$ 的\Term{perfect-dimension}{完美维数} $\leq n$，并写作 $\operatorname{parf.dim}(F)\leq n$。使 $\operatorname{parf.dim}(F)\leq n$ 成立的最小整数 $n$，称为 $F$ 的\Term{perfect-dimension}{完美维数}（若不存在这样的整数，则称 $F$ 的\Term{perfect-dimension}{完美维数}为无穷）。
\end{SGAresult}
当 $\mathscr S$ 为\Term{point-site}{点景}，且 $\mathcal C_0$ 为 $\mathcal C$ 中由\Term{projective-object}{投射}\Term{object}{对象}构成的\Term{full-subcategory}{满子范畴}时，\Term{perfect-dimension}{完美维数}就是通常的\Term{projective-dimension}{投射维数}。

\begin{SGAresult}{命题}{4.15}
设 $a,b,n\in\mathbb Z$，$a\leq b$、$n\geq0$，$F\in\operatorname{ob}D(\mathcal C)$，且当 $i\notin[a,b]$ 时 $F^i=0$（相应地，$H^i(F)=0$）。若对每个 $i\in[a,b]$，$F^i$（相应地，$H^i(F)$）的\Term{perfect-dimension}{完美维数} $\leq n+(i-a)$，则 $F$ 的\Term{perfect-amplitude}{完美振幅}包含于 $[a-n,b]$。因此，若每个 $i\in[a,b]$ 对应的 $F^i$（相应地，$H^i(F)$）都\Term{perfect-complex}{完美}，则 $F$ 也\Term{perfect-complex}{完美}。
\end{SGAresult}
\SourcePage{138}
\begin{proof}固定 $a,n$，对 $b$ 归纳；\Term{truncation}{截断} $F$ 并应用（4.11）（参看（2.11）的证明）。\end{proof}

\begin{SGAresult}{\Term{lemma}{引理}}{4.16}
设 $a,b\in\mathbb Z$，$a\leq b$，$F\in\operatorname{ob}D(\mathcal C)$，且 $i\notin[a,b]$ 时 $F^i=0$。假定 $\operatorname{parf.amp}(F)\subset[a,b]$，且 $i>a$ 时 $F^i\in\operatorname{ob}\mathcal C_0$。则 $F^a$ \Term{local-scheme}{局部}为 $\mathcal C_0$ 的\Term{object}{对象}。
\end{SGAresult}
\begin{proof}
\Term{truncation}{截断} $F$ 得到\Term{exact-sequence}{正合列}
\[0\longrightarrow F'\longrightarrow F\longrightarrow F^a[-a]\longrightarrow0,\]
其中 $\operatorname{parf.amp}(F')\subset[a+1,b+1]$。由（4.11）与（4.13）可得 $\operatorname{parf.dim}(F^a)=0$，即 $F^a$ \Term{local-scheme}{局部}为 $\mathcal C_0$ 的\Term{object}{对象}。
\end{proof}

\begin{SGAresult}{命题}{4.17}
设 $F',F''\in\operatorname{ob}D(\mathcal C_S)$，$a,b\in\mathbb Z$，$a\leq b$。考虑下列条件：
\begin{enumerate}
\item[(i)] $\operatorname{parf.amp}(F')\subset[a,b]$ 且 $\operatorname{parf.amp}(F'')\subset[a,b]$；
\item[(ii)] $\operatorname{parf.amp}(F'+F'')\subset[a,b]$。
\end{enumerate}
有（i）$\Rightarrow$（ii）。若 $\mathcal C_0$ \Term{local-scheme}{局部}地对\Term{direct-summand}{直和因子}\Term{stable}{稳定}（即对 $E',E''\in\operatorname{ob}\mathcal C_X$、$X\in\operatorname{ob}\mathscr S$，关系 $E'+E''\in\operatorname{ob}\mathcal C_0$ 蕴含 $E',E''$ \Term{local-scheme}{局部}为 $\mathcal C_0$ 的\Term{object}{对象}），则（i）$\Leftrightarrow$（ii）。
\end{SGAresult}
\begin{proof}
蕴含（i）$\Rightarrow$（ii）是（4.11）的一个特例。证明逆向蕴含。置 $F=F'+F''$。因为 $F$ 为\Term{pseudo-coherent}{伪凝聚}的，$F',F''$ 也如此（2.12）。另一方面，$i\notin[a,b]$ 时 $H^i(F)=0$ 蕴含 $H^i(F')=H^i(F'')=0$。由（2.10 a）），\Term{localization}{局部化}后可设 $i\geq a+1$ 时 $F'^i,F''^i\in\operatorname{ob}\mathcal C_0$，且 $i>b$ 时 $F'^i=F''^i=0$；再作\Term{truncation}{截断}，可令
\SourcePage{139}
$i<a$ 时 $F'^i=F''^i=0$。因为 $\operatorname{parf.dim}(F)\subset[a,b]$，由（4.16），$F^a$ \Term{local-scheme}{局部}为 $\mathcal C_0$ 的\Term{object}{对象}。而 $F^a=F'^a+F''^a$，所以由 $\mathcal C_0$ \Term{local-scheme}{局部}地对\Term{direct-summand}{直和因子}\Term{stable}{稳定}，$F'^a,F''^a$ 也\Term{local-scheme}{局部}为 $\mathcal C_0$ 的\Term{object}{对象}。
\end{proof}

\begin{SGAresult}{命题}{4.18}
设三元组 $(\mathscr S',\mathcal C',\mathcal C'_0)$ 满足与 $(\mathscr S,\mathcal C,\mathcal C_0)$ 相同的假设（4.0）。设 $f:\mathscr S'\to\mathscr S$ 为\Term{site}{景}的\Term{morphism}{态射}（由 $f^{-1}:\mathscr S\to\mathscr S'$ 定义），$u:D^*(\mathcal C)\to D^*(\mathcal C')$（$*=-,+,\varnothing$）为 $f$ 上方的\Term{functor}{函子}（参看（2.13））。若 $u$ 将 $\mathcal C_0$ 的\Term{object}{对象}变为 $\mathcal C'_0$-\Term{perfect-complex}{完美复形}，则 $u$ 定义\Term{functor}{函子}
\[D^*(\mathcal C)_{\mathcal C_0\text{-parf}}\longrightarrow D^*(\mathcal C')_{\mathcal C'_0\text{-parf}}.\]
\end{SGAresult}
\begin{proof}
设 $F\in\operatorname{ob}D^*(\mathcal C)_{\mathcal C_0\text{-parf}}$，需证明 $u(F)\in\operatorname{ob}D^*(\mathcal C')_{\mathcal C'_0\text{-parf}}$。问题是\Term{local-scheme}{局部}的，故可设 $F$ 严格\Term{perfect-complex}{完美}。只须逐步\Term{resolution}{分解}，并注意 $D^*(\mathcal C')_{\mathcal C'_0\text{-parf}}$ 是 $D^*(\mathcal C')$ 的\Term{triangulated-subcategory}{三角子范畴}（4.10）。
\end{proof}

\textbf{注记 4.18.1。}假定存在 $m,n\in\mathbb Z$（$m\leq n$），使 $F\in\operatorname{ob}\mathcal C_0$ 蕴含 $\mathcal C'_0\text{-}\operatorname{parf.amp}(u(F))\subset[m,n]$。则对 $F\in\operatorname{ob}D^*(\mathcal C)$，有
\[\mathcal C_0\text{-}\operatorname{parf.amp}(F)\subset[a,b]
\ \Longrightarrow\ 
\mathcal C'_0\text{-}\operatorname{parf.amp}(u(F))\subset[a+m,b+n].\]
证明留给读者作练习（使用（4.11））。

\begin{SGAresult}{推论}{4.19}
在（2.16）的假设下，若 $E'\in\operatorname{ob}D^-({}_{A'}\mathscr S'_A)$ 作为 $D^-({}_{A'}\mathscr S')$ 的\Term{object}{对象}是\Term{perfect-complex}{完美}的，且 $E\in\operatorname{ob}D^-({}_A\mathscr S)$ 是\Term{perfect-complex}{完美}的，则 $E'\otimes_A^{\mathrm L}E$ \Term{perfect-complex}{完美}。
\end{SGAresult}
\begin{proof}
可设 $E'\in\operatorname{ob}D^-({}_{A'}S'_A)$（$S'$ 是 $\mathscr S'$ 的\Term{final-object}{终对象}），将（4.18）应用于 $u=E'\otimes_A^{\mathrm L}(\,\cdot\,):D^-({}_A\mathscr S)\to D^-({}_{A'}\mathscr S')$ 即可。
\end{proof}
特别地：
\SourcePage{140}
\begin{SGAresult}{推论}{4.19.1}
a）\Term{functor}{函子} $\mathrm Lf^*:D^-({}_A\mathscr S)\to D^-({}_{A'}\mathscr S')$ 诱导\Term{functor}{函子}
\[\mathrm Lf^*:D^-({}_A\mathscr S)_{\mathrm{parf}}\longrightarrow D^-({}_{A'}\mathscr S')_{\mathrm{parf}}.\]
b）若 $A$ 交换，则\Term{functor}{函子} $\otimes_A^{\mathrm L}:D^-(\mathscr S)\times_{\mathscr S}D^-(\mathscr S)\to D^-(\mathscr S)$ 诱导\Term{functor}{函子}
\[\otimes_A^{\mathrm L}:D^-(\mathscr S)_{\mathrm{parf}}\times_{\mathscr S}D^-(\mathscr S)_{\mathrm{parf}}\longrightarrow D^-(\mathscr S)_{\mathrm{parf}}.\]
\end{SGAresult}

\textbf{注记 4.19.2。}可以用\Term{perfect-amplitude}{完美振幅}更精确地表述（4.19）：
\[
\begin{gathered}
\operatorname{parf.amp}(E)\subset[a,b],\quad\operatorname{parf.amp}(E')\subset[a',b']\\
\Longrightarrow\quad\operatorname{parf.amp}(E'\otimes_A^{\mathrm L}E)\subset[a+a',b+b'].
\end{gathered}
\]
特别地，
\[\operatorname{parf.amp}(E)\subset[a,b]\ \Longrightarrow\ \operatorname{parf.amp}(\mathrm Lf^*E)\subset[a,b].\]

\SGAsection{5}{有限 Tor 维数与完美性}
\textbf{5.0。}读者想必已经提出这样一个问题：一个\Term{pseudo-coherent}{伪凝聚}\Term{complex}{复形} $E$ 还缺什么才会\Term{perfect-complex}{完美}？例如，仅有\Term{cohomology}{上同调}\Term{bounded}{有界}是否足够？答案是否定的。取 $\mathcal C$ 为交换 Noether 环 $A$ 上的\Term{module-category}{模范畴}，$\mathcal C_0$ 为\Term{finite-type}{有限型}\Term{projective-object}{投射} $A$-模组成的\Term{full-subcategory}{满子范畴}。任意\Term{finite-type}{有限型} $A$-模 $E$ 都是\Term{pseudo-coherent}{伪凝聚}且\Term{cohomology}{上同调}\Term{bounded}{有界}的，但 $E$ \Term{perfect-complex}{完美}还要求它的\Term{projective-dimension}{投射维数}\Term{finite}{有限}（\Term{equivalence}{等价}地，\Term{tor-dimension}{Tor 维数}\Term{finite}{有限}）。这个例子提示：\Term{pseudo-coherent}{伪凝聚}\Term{complex}{复形}与\Term{perfect-complex}{完美复形}的差别在于\Term{cohomological-dimension}{上同调维数}。现在我们要在
\SourcePage{141}
\Term{ringed-topos}{赋环拓扑斯}的\Term{module-category}{模范畴}中考察这个问题（参看（3.3））。先要讨论\Term{finite}{有限}\Term{tor-dimension}{Tor 维数}的一些预备性质。由于文献已有论述，这里将从略：参看（[H] II 4）及（SGA 4 XVII 4.1.9）。

\begin{SGAresult}{命题}{5.1}
设 $(\mathscr S,A)$ 为\Term{ringed-topos}{赋环拓扑斯}，$X\in\operatorname{ob}\mathscr S$，$F\in\operatorname{ob}D^-({}_AX)$，$m,n\in\mathbb Z$，$m\leq n$。下列条件\Term{equivalence}{等价}：

（i）$F$ 在 $D^-({}_AX)$ 中\Term{isomorphism}{同构}于一个\Term{complex}{复形} $F'$，使得每个 $F'^i$ 都\Term{flat}{平坦}，且 $i\notin[m,n]$ 时为零；

（ii）对每个右 $A_X$-模 $E$，当 $i\notin[m,n]$ 时都有 $H^i(E\otimes_A^{\mathrm L}F)=0$。
\end{SGAresult}
\begin{proof}
（i）$\Rightarrow$（ii）显然；（ii）$\Rightarrow$（i）可通过\Term{truncation}{截断}由下列\Term{lemma}{引理}得到。
\end{proof}

\begin{SGAresult}{\Term{lemma}{引理}}{5.1.1}
（参看（4.16）。）设 $F\in\operatorname{ob}D^-({}_AX)$。假定 $F$ 满足（5.1）的条件（ii），$i\notin[m,n]$ 时 $F^i=0$，且 $i\ne m$ 时 $F^i$ \Term{flat}{平坦}。则 $F^m$ \Term{flat}{平坦}。
\end{SGAresult}
\begin{proof}
有\Term{distinguished-triangle}{可区别三角形}
\[F'\longrightarrow F\longrightarrow F^m[-m]\xrightarrow{+1}F',\]
其中 $F'$ 各项\Term{flat}{平坦}，且在区间 $[m+1,n]$ 外为零（设 $m<n$）。对任意右 $A_X$-模 $E$，应用\Term{functor}{函子} $E\otimes_A^{\mathrm L}(\,\cdot\,)$ 并写出\Term{cohomology}{上同调}\Term{exact-sequence}{正合列}即可。
\end{proof}

\SourcePage{142}
\begin{SGAresult}{定义}{5.2}
若 $D^-({}_A\mathscr S)$ 的\Term{object}{对象} $F$ 满足（5.1）的\Term{equivalence}{等价}条件，则称 $F$ 的\Term{flat-amplitude}{平坦振幅}包含于区间 $[m,n]$，并记作
\[\operatorname{tor.amp}(F)\subset[m,n].\]
若存在整数 $m,n$ 使上述条件成立，则称 $F$ 的\Term{flat-amplitude}{平坦振幅}\Term{finite}{有限}（或\Term{tor-dimension}{Tor 维数}\Term{finite}{有限}）。
\end{SGAresult}

\textbf{5.3。}设 $X\in\operatorname{ob}\mathscr S$。以 $D^-({}_AX)_{\mathrm{torf}}$ 表示 $D^-({}_AX)$ 中由\Term{finite}{有限}\Term{tor-dimension}{Tor 维数}\Term{object}{对象}组成的\Term{full-subcategory}{满子范畴}。由（5.1）的判据（ii），这是 $D^-({}_AX)$ 的\Term{triangulated-subcategory}{三角子范畴}。当 $X$ 变化时，\Term{category}{范畴} $D^-({}_AX)_{\mathrm{torf}}$ 构成 $D^-({}_A\mathscr S)$ 的一个\Term{triangulated-subcategory}{三角}子 $\mathscr S$-\Term{category}{范畴}（1.1.4），记作 $D^-({}_A\mathscr S)_{\mathrm{torf}}$。\Term{functor}{函子} $\otimes_A^{\mathrm L}$ 诱导\Term{exact}{正合}的 $\mathscr S$-\Term{functor}{函子}
\[\otimes_A^{\mathrm L}:D^b(\mathscr S_A)\times_{\mathscr S}D^-({}_A\mathscr S)_{\mathrm{torf}}\longrightarrow D^b({}_{\mathbb Z}\mathscr S).\]

\textbf{5.4。}左 $A$-模 $F$ 的\Term{flat-amplitude}{平坦振幅}包含于 $[-n,0]$（$n\in\mathbb N$），当且仅当 $F$ 在通常意义下的\Term{tor-dimension}{Tor 维数} $\leq n$，即 $F$ 有\Term{length}{长度}为 $n$、由\Term{flat}{平坦} $A$-模构成的左\Term{resolution}{消解}。

\SourcePage{143}
\textbf{5.5。}设 $X\in\operatorname{ob}\mathscr S$，$m,n\in\mathbb Z$，$m\leq n$。$F\in\operatorname{ob}D^-({}_AX)$ 的\Term{flat-amplitude}{平坦振幅}包含于 $[m,n]$，当且仅当存在\Term{covering-family}{覆盖族} $(X_i\to X)$，使 $F|X_i$ 的\Term{flat-amplitude}{平坦振幅}包含于 $[m,n]$。特别地，若 $\mathscr S_{/X}$ 为\Term{finite-type}{有限型}（SGA 4 VI 1.7），且 $F$ \Term{local-scheme}{局部}为\Term{finite}{有限}\Term{tor-dimension}{Tor 维数}，则 $F$ 的\Term{tor-dimension}{Tor 维数}\Term{finite}{有限}。

\begin{SGAresult}{命题}{5.6}
在（2.16）的假设下，若 $E'$ 作为 $D^-({}_{A'}\mathscr S')$ 的\Term{object}{对象}，其\Term{flat-amplitude}{平坦振幅}包含于 $[m',n']$，且 $E$ 的\Term{flat-amplitude}{平坦振幅}包含于 $[m,n]$，则 $E'\otimes_A^{\mathrm L}E$ 的\Term{flat-amplitude}{平坦振幅}包含于 $[m+m',n+n']$。
\end{SGAresult}
\begin{proof}
这是直接的：使用\Term{finite}{有限}\Term{tor-dimension}{Tor 维数}的（5.1（i））\Term{form}{形式}，并注意，若 $E$ 为\Term{flat}{平坦}左 $A$-模，则 $f^{-1}(E)$ 为\Term{flat}{平坦}左 $f^{-1}(A)$-模（SGA 4 IV 附录 II）。
\end{proof}
特别地：
\begin{SGAresult}{推论}{5.6.1}
a）\Term{functor}{函子} $\mathrm Lf^*:D^-({}_A\mathscr S)\to D^-({}_{A'}\mathscr S')$ 诱导\Term{functor}{函子}
\[\mathrm Lf^*:D({}_A\mathscr S)_{\mathrm{torf}}\longrightarrow D({}_{A'}\mathscr S')_{\mathrm{torf}}.\]
b）若 $A$ 交换，则\Term{functor}{函子} $\otimes_A^{\mathrm L}:D^-(\mathscr S)\times_{\mathscr S}D^-(\mathscr S)\to D^-(\mathscr S)$ 诱导\Term{functor}{函子}
\[\otimes_A^{\mathrm L}:D(\mathscr S)_{\mathrm{torf}}\times_{\mathscr S}D(\mathscr S)_{\mathrm{torf}}\longrightarrow D(\mathscr S)_{\mathrm{torf}}.\]
\end{SGAresult}

\textbf{注记 5.6.2。}设 $E\in\operatorname{ob}D^-({}_AS)$（$S$ 为 $\mathscr S$ 的\Term{final-object}{终对象}），$m,n\in\mathbb Z$，$m\leq n$。若 $E$ 的\Term{flat-amplitude}{平坦振幅}包含于 $[m,n]$，则由（5.6），对 $\mathscr S$ 的每个点 $s$，$E_s$ 也如此。若 $\mathscr S$ 有\Term{enough-points}{足够多的点}，则逆命题成立（使用\Term{flat-amplitude}{平坦振幅}的（5.1（ii））\Term{form}{形式}，以及 $\otimes_A^{\mathrm L}$ 与取\Term{fiber}{纤维}
\SourcePage{144}
交换这一事实）。

\textbf{练习 5.7。}a）假定 $A$ 交换，且 $(\mathscr S',A')$ 在 $(\mathscr S,A)$ 上\Term{faithfully-flat}{忠实平坦}（参看（2.16.2））。$E\in\operatorname{ob}D^-({}_AS)$\footnote{$S$ 为 $\mathscr S$ 的\Term{final-object}{终对象}。} 的\Term{flat-amplitude}{平坦振幅}包含于 $[m,n]$，当且仅当 $f^*(E)$ 如此。

b）设 $F\in\operatorname{ob}D^+({}_AS)$\footnotemark[\value{footnote}]。证明下列条件\Term{equivalence}{等价}：
\begin{enumerate}
\item[(i)] $F$ 在 $D^+({}_AS)$ 中\Term{isomorphism}{同构}于一个\Term{complex}{复形} $F'$，使每个 $F'^i$ 都\Term{injective}{内射}，且 $i\notin[m,n]$ 时为零；
\item[(ii)] 对每个左 $A_S$-模 $E$，$i\notin[m,n]$ 时 $\mathcal Ext^i(E,F)=0$。
\end{enumerate}
若 $F$ 满足上述条件，则称 $F$ 的\Term{injective-amplitude}{内射振幅}包含于 $[m,n]$。$D^+({}_AS)$ 中由\Term{finite}{有限}\Term{injective-amplitude}{内射振幅}\Term{object}{对象}组成的\Term{full-subcategory}{满子范畴}是\Term{triangulated-subcategory}{三角子范畴}，记作 $D^+({}_AS)_{\mathrm{injf}}$。证明：若 $A$ 交换，则\Term{functor}{函子}
\[\operatorname{R\mathcal Hom}:D(\mathscr S)^\circ\times D^+(\mathscr S)\longrightarrow D(\mathscr S)\]
诱导\Term{functor}{函子}
\[\operatorname{R\mathcal Hom}:D(\mathscr S)_{\mathrm{torf}}^\circ\times D(\mathscr S)_{\mathrm{injf}}\longrightarrow D(\mathscr S)_{\mathrm{injf}}.\]

现在可以回答（5.0）的问题，并陈述本节的主要结果。
\begin{SGAresult}{\Term{theorem}{定理}}{5.8}
设 $(\mathscr S,A)$ 为\Term{ringed-topos}{赋环拓扑斯}，$E\in\operatorname{ob}D({}_A\mathscr S)$，$m,n\in\mathbb Z$，$m\leq n$。下列条件\Term{equivalence}{等价}：
\begin{enumerate}
\item[(i)] $\operatorname{parf.amp}(E)\subset[m,n]$（4.8）；
\item[(ii)] $E$ 为 $(m-1)$-\Term{pseudo-coherent}{伪凝聚}的（2.15），且 $\operatorname{tor.amp}(E)\subset[m,n]$（5.2）。
\end{enumerate}
\end{SGAresult}
由于“\Term{perfect-complex}{完美}”显然蕴含“\Term{pseudo-coherent}{伪凝聚}”，可得：
\SourcePage{145}
\begin{SGAresult}{推论}{5.8.1}
$E$ \Term{perfect-complex}{完美}（4.8），当且仅当 $E$ 为\Term{pseudo-coherent}{伪凝聚}的（2.15），且\Term{local-scheme}{局部}具有\Term{finite}{有限}\Term{tor-dimension}{Tor 维数}（5.2）。
\end{SGAresult}
证明需要几个\Term{lemma}{引理}。

\begin{SGAresult}{\Term{lemma}{引理}}{5.8.2}
设 $B$ 为 $\mathscr S$ 中的环。若 $E$ 为左 $A$-模，$G$ 为右 $B$-模，$F$ 为 $A$-$B$-\Term{bimodule}{双模}（左 $A$、右 $B$），则存在\Term{canonical}{典范}的\Term{functorial}{函子性}\Term{homomorphism}{同态}
\[\mathcal Hom_B(F,G)\otimes_A E\longrightarrow
\mathcal Hom_B(\mathcal Hom_A(E,F),G).\tag{5.8.2.1}\]
当 $G$ \Term{injective}{内射}且 $E$ 为\Term{finite-presentation}{有限表示}时，此\Term{homomorphism}{同态}是\Term{isomorphism}{同构}。
\end{SGAresult}
\begin{proof}
先定义\Term{arrow}{箭头}。由 Cartan 所钟爱的\Term{adjunction}{伴随}公式，这\Term{equivalence}{等价}于定义一个 $B$-\Term{module-homomorphism}{模同态}
\[\mathcal Hom_B(F,G)\otimes_A E\otimes_{\mathbb Z_S}\mathcal Hom_A(E,F)\longrightarrow G.\tag{*}\]
\footnote{以 $S$ 表示 $\mathscr S$ 的\Term{final-object}{终对象}，$\mathbb Z_S$ 表示值为 $\mathbb Z$ 的\Term{constant-sheaf}{常值层}。}
而有\Term{canonical}{典范}的 $A$-$B$-\Term{bimodule}{双模}\Term{homomorphism}{同态}
\[E\otimes_{\mathbb Z_S}\mathcal Hom_A(E,F)\longrightarrow F\tag{**}\]
（它由 $\mathcal Hom_A(E,F)$ 的恒等\Term{arrow}{箭头}通过\Term{adjunction}{伴随}得到）。同样，有\Term{canonical}{典范}的 $B$-\Term{module-homomorphism}{模同态}
\[\mathcal Hom_B(F,G)\otimes_A F\longrightarrow G.\tag{***}\]
于是以（**）和（***）的复合定义（*）。要看出 $G$ \Term{injective}{内射}、$E$ 为\Term{finite-presentation}{有限表示}时（5.8.2.1）是\Term{isomorphism}{同构}，只须注意：当 $G$ \Term{injective}{内射}时，此\Term{arrow}{箭头}的源与目标都是关于 $E$ 的右\Term{exact}{正合}\Term{functor}{函子}，且在 $E=A$ 时该\Term{arrow}{箭头}诱导\Term{identity-morphism}{恒等态射}。
\end{proof}

\SourcePage{146}
下面是（5.8）的一个特例。
\begin{SGAresult}{\Term{lemma}{引理}}{5.8.3}
设 $E$ 为左 $A$-模。下列条件\Term{equivalence}{等价}：
\begin{enumerate}
\item[(i)] $E$ \Term{local-scheme}{局部}为\Term{finite-free}{有限型自由模}的\Term{direct-summand}{直和因子}；
\item[(ii)] $E$ \Term{flat}{平坦}且为\Term{finite-presentation}{有限表示}。
\end{enumerate}
\end{SGAresult}
\begin{proof}[证明（依据 BOURBAKI《\Term{commutative-algebra}{交换代数}》第 I 章，\S 2，练习 15）]
（i）显然蕴含（ii）。证明（ii）$\Rightarrow$（i）。由（1.3.1），只须证明\Term{functor}{函子} $\mathcal Hom_A(E,\,\cdot\,)$ \Term{exact}{正合}。设 $F\to F''\to0$ 为（左）$A$-模的\Term{exact-sequence}{正合列}。需证明序列
\[\mathcal Hom_A(E,F)\longrightarrow\mathcal Hom_A(E,F'')\longrightarrow0\]
\Term{exact}{正合}。只须证明，对任意\Term{injective}{内射} $\mathbb Z_S$-模 $I$，应用\Term{functor}{函子} $\mathcal Hom_{\mathbb Z}(\,\cdot\,,I)$ 所得的序列\Term{exact}{正合}。有如下\Term{commutative-diagram}{交换图}，其中竖直\Term{arrow}{箭头}为（3.17.1）中的\Term{canonical}{典范}\Term{arrow}{箭头}：
\[\begin{tikzcd}[column sep=small]
0\arrow[r]&\mathcal Hom_{\mathbb Z_S}(F'',I)\otimes_A E\arrow[r]\arrow[d]&\mathcal Hom_{\mathbb Z_S}(F,I)\otimes_A E\arrow[d]\\
0\arrow[r]&\mathcal Hom_{\mathbb Z_S}(\mathcal Hom_A(E,F''),I)\arrow[r]&\mathcal Hom_{\mathbb Z_S}(\mathcal Hom_A(E,F),I).
\end{tikzcd}\]
由（5.8.2），竖直\Term{arrow}{箭头}均为\Term{isomorphism}{同构}。由于 $E$ \Term{flat}{平坦}，上行\Term{exact}{正合}，所以下行\Term{exact}{正合}，证毕。
\end{proof}

\begin{proof}[（5.8）的证明]
蕴含（i）$\Rightarrow$（ii）显然由定义得到。证明（ii）$\Rightarrow$（i）。问题是\Term{local-scheme}{局部}的，故可设 $E$ 严格 $(m-1)$-\Term{pseudo-coherent}{伪凝聚}。另一方面，$E$ 在\Term{degree}{次数} $<m$ \Term{acyclic}{无上同调}，故\Term{truncation}{截断}\Term{arrow}{箭头}
\SourcePage{147}
\[\begin{tikzcd}
\cdots\arrow[r]&E^{m-1}\arrow[r]\arrow[d]&E^m\arrow[r]\arrow[d]&E^{m+1}\arrow[r]\arrow[d,equal]&\cdots\\
\cdots\arrow[r]&0\arrow[r]&E^m/B^m\arrow[r]&E^{m+1}\arrow[r]&\cdots
\end{tikzcd}\]
是\Term{quasi-isomorphism}{拟同构}。以 $E'$ 表示\Term{truncation}{截断}后的\Term{complex}{复形}。由构造，$E'$ 的\Term{flat-amplitude}{平坦振幅}包含于 $[m,n]$，且 $i>m$ 时 $E'^i$ 为\Term{finite-free}{有限型自由模}。故由（5.1.1），$E'^m$ \Term{flat}{平坦}。另一方面，有\Term{distinguished-triangle}{可区别三角形}
\[E''\longrightarrow E'\longrightarrow E'^m[-m]\xrightarrow{+1}E'',\]
其中 $E''$ 严格\Term{perfect-complex}{完美}。由（2.6），$E'^m[-m]$ 为 $(m-1)$-\Term{pseudo-coherent}{伪凝聚}，即由（2.9），$E'^m$ 为\Term{finite-presentation}{有限表示}。由（5.8.3），$E'^m$ \Term{local-scheme}{局部}为\Term{finite-free}{有限型自由模}的\Term{direct-summand}{直和因子}，从而 $\operatorname{parf.amp}(E)\subset[m,n]$。
\end{proof}

\textbf{注记 5.8.4。}a）利用（2.5），可重新得到：\Term{perfect-complex}{完美复形}组成 $D({}_A\mathscr S)$ 的一个\Term{triangulated-subcategory}{三角子范畴}。

b）\Term{perfect-complex}{完美复形}即使\Term{cohomology}{上同调}\Term{bounded}{有界}，也不一定具有\Term{finite}{有限}\Term{tor-dimension}{Tor 维数}。但若 $\mathscr S$ 为\Term{finite-type}{有限型}，则确实如此（5.5）。以 $D({}_A\mathscr S)_{\underline{\mathrm{parf}}}$ 表示 $D^b({}_A\mathscr S)$ 中由\Term{perfect-complex}{完美}且\Term{tor-dimension}{Tor 维数}\Term{finite}{有限}的\Term{complex}{复形}构成的\Term{full-subcategory}{满}子 $\mathscr S$-\Term{category}{范畴}。由（2.16.1 b））及（5.3）立即得到：若 $A$ 交换，则\Term{functor}{函子} $\otimes_A^{\mathrm L}$ 诱导\Term{functor}{函子}
\[\otimes_A^{\mathrm L}:D(\mathscr S)_{\underline{\mathrm{parf}}}\times_{\mathscr S}D^b(\mathscr S)_{\mathrm{coh}}\longrightarrow D^b(\mathscr S)_{\mathrm{coh}}.\]
\SourcePage{148}
这个公式是\Term{scheme}{概形}上\Term{intersection-theory}{相交理论}的起点。另一方面应当注意，两个\Term{cohomology}{上同调}\Term{bounded}{有界}的\Term{pseudo-coherent}{伪凝聚}\Term{complex}{复形}的（\Term{derive}{导出}）\Term{tensor-product}{张量积}，其\Term{cohomology}{上同调}\Term{generic}{一般}并不\Term{bounded}{有界}。

\begin{SGAresult}{推论}{5.9}
假定 $\mathscr S=(\mathrm{Ens})$（因而 $A$ “就是”一个通常的环），且每个左 $A$-模都具有\Term{finite}{有限}\Term{tor-dimension}{Tor 维数}（例如 $A$ 的全局\Term{cohomological-dimension}{上同调维数}\Term{finite}{有限}时即是如此，见 MULT VI 2.6 及 EGA $0_{\mathrm{IV}}$ 17.2.8）。则
\begin{align*}
\text{（i）}\quad&D^b({}_A\mathscr S)=D({}_A\mathscr S)_{\mathrm{torf}},\\
\text{（ii）}\quad&D^b({}_A\mathscr S)_{\mathrm{coh}}=D({}_A\mathscr S)_{\underline{\mathrm{parf}}}.
\end{align*}
换言之，$E\in\operatorname{ob}D({}_A\mathscr S)$ 的\Term{tor-dimension}{Tor 维数}\Term{finite}{有限}（相应地，\Term{tor-dimension}{Tor 维数}\Term{finite}{有限}且\Term{perfect-complex}{完美}），当且仅当 $E$ 的\Term{cohomology}{上同调}\Term{bounded}{有界}（相应地，\Term{cohomology}{上同调}\Term{bounded}{有界}且\Term{pseudo-coherent}{伪凝聚}）。
\end{SGAresult}
\begin{proof}
由（5.8），只须证明（i）。设 $E\in\operatorname{ob}D^b({}_A\mathscr S)$。作\Term{extension-by-zero}{零延拓}后可设 $E\in\operatorname{ob}D^b({}_AS)$，其中 $S$ 为 $\mathscr S$ 的\Term{final-object}{终对象}。对 $E$ 的非零\Term{cohomology}{上同调}\Term{object}{对象}的数目归纳，通过逐步\Term{resolution}{分解}，归结为 $E$ 是普通 $A$-模的情形，结论即得。
\end{proof}
更\Term{generic}{一般}地，可以陈述如下推论。

\begin{SGAresult}{推论}{5.10}
设 $X\in\operatorname{ob}\mathscr S$。假定 $\mathscr S_{/X}$ 有\Term{enough-points}{足够多的点}（SGA 4 IV 6.4.1），且对 $\mathscr S_{/X}$ 的每个点 $s$，每个左 $A_s$-模都具有\Term{finite}{有限}\Term{tor-dimension}{Tor 维数}。则
\[D^b({}_AX)_{\mathrm{coh}}=D^b({}_AX)_{\mathrm{parf}}.\]
\SourcePage{149}
换言之，$E\in\operatorname{ob}D^b({}_AX)$ 为\Term{pseudo-coherent}{伪凝聚}的，当且仅当 $E$ \Term{perfect-complex}{完美}。
\end{SGAresult}
\begin{proof}
设 $E\in\operatorname{ob}D^b({}_AX)_{\mathrm{coh}}$，$s$ 为 $\mathscr S_{/X}$ 的点。需证明 $E$ “在 $s$ 的\Term{neighborhood}{邻域}上”\Term{perfect-complex}{完美}，即存在 $\mathscr S$ 中位于 $X$ 上方的 $s$-\Term{pointed-object}{带点对象} $U$，使 $E|U$ \Term{perfect-complex}{完美}。因为\Term{functor}{函子} $s^*$ \Term{exact}{正合}，由（2.16.1 a））有 $E_s\in\operatorname{ob}D^b(A_s)_{\mathrm{coh}}$。故由（5.9），$E_s$ \Term{perfect-complex}{完美}。于是断言由（3.7.2）及下列\Term{lemma}{引理}得到。
\end{proof}

\begin{SGAresult}{\Term{lemma}{引理}}{5.10.1}
设 $s$ 为 $\mathscr S$ 的点，$F\in\operatorname{ob}D(A_s)_{\mathrm{parf}}$。存在 $\mathscr S$ 中的 $s$-\Term{pointed-object}{带点对象} $U$ 及 $F'\in\operatorname{ob}D({}_AU)_{\mathrm{parf}}$，使 $F'_s=F$。
\end{SGAresult}
\begin{proof}
当 $F$ 为\Term{finite-free}{有限型自由} $A_s$-模时，这是显然的。再设 $F$ 是\Term{finite-type}{有限型}\Term{projective-object}{投射}模，因而是\Term{projector}{投影算子} $p:L\to L$ 的像，其中 $L$ 为\Term{finite-free}{有限型自由模}。存在 $s$ 的一个“\Term{neighborhood}{邻域}” $U$、一个自由 $A_U$-模 $L'$ 以及\Term{homomorphism}{同态} $p':L'\to L'$，使 $p'_s=p$。由于 $(p'_s)^2=p'_s$，缩小 $U$ 后可设 $p'^2=p'$。故 $F'=\operatorname{Im}(p')$ 是 $L'$ 的\Term{direct-summand}{直和因子}，并满足 $F'_s=F$。立刻可推广到 $F$ 严格\Term{perfect-complex}{完美}的情形，即 $F$ 为由\Term{finite-type}{有限型}\Term{projective-object}{投射}模组成、\Term{degree}{次数}\Term{bounded}{有界}的\Term{complex}{复形}。
\end{proof}

\textbf{例 5.11。}推论（5.10）特别适用于 $(\mathscr S,A)$ 由一个取值为\Term{regular-local-ring}{正则局部环}的\Term{ringed-space}{赋环空间}所定义的情形，例如\Term{regular-scheme}{正则概形}，或 $\mathbb C$ 上的\Term{smooth}{光滑}\Term{complex-analytic-space}{复解析空间}。
